\section*{Conventions and scope}
We fix Grothendieck universes $\U\in\V$ with $\N\in\U$; a category called small has its sets of objects and arrows in $\U$. Categories of all $\U$-small objects and their presheaves are formed in $\V$. For a small category $X$, $NX$ is its nerve and $BX=|NX|$ its classifying space. All homotopy classes are unpointed unless a subscript $*$ occurs. We put $\Hom_{\Cathot}(X,K)=\pi_0\Fun(X,K)$, using finite zigzags of natural transformations. We write $\W_\infty$ for functors whose nerves are weak homotopy equivalences and $\Hot=\Cat[\W_\infty^{-1}]$. Equality, categorical isomorphism, categorical equivalence, isomorphism in $\Cathot$, and isomorphism in $\Hot$ are distinct notions.

For $u:A\to B$ and $b\in B$, $A/b$ has objects $(a,u(a)\to b)$ and $b\backslash A$ has objects $(a,b\to u(a))$. A functor is final in the covariant convention when each $b\backslash A$ is nonempty and connected. Presheaf coefficients have variance $X^{\op}\to D(e)$. A sieve is full and closed under sources of arrows, and a cosieve is full and closed under targets. We translate the manuscript's terminology according to this convention at each use.

A fundamental localizer $\W$ is used with the following precise version of Loc(1)--(3). Loc(1): $\W$ contains identities, has two-out-of-three, and if $i:A\to B$, $r:B\to A$ satisfy $ri=\id_A$ and $ir\in\W$, then $i,r\in\W$. Loc(2): if $A$ has a terminal object then $A\to e$ belongs to $\W$. Loc(3), relative Theorem A: for a strictly commutative triangle $A\xrightarrow{u}B\xrightarrow{v}C$, if $u/c:A/c\to B/c$ belongs to $\W$ for every $c\in C$, then $u\in\W$. Here the slice structures are those induced by $vu$ and $v$. We call $A$ $\W$-aspherical when $A\to e$ belongs to $\W$. Our homotopy-theoretic conclusions use $\W_\infty$; formal localization statements explicitly indicate when another class is allowed.

A finite category has finitely many arrows, including identities. A finite homotopy type, when this expression is used in the compact-generation argument, means a homotopy type represented by a finite CW complex; it does not mean a space with finitely many homotopy groups or a nerve of a finite category. Empty categories and empty spaces are permitted.
